\section{unknot と Hopf 絡み目}\label{sec:hopf}

\subsection{virtual 正規化}
probe 側の基底 $K_{(n)}(y; t, q; 1, -1, q^{1/2}, -q^{1/2})$（monic な Askey--Wilson 多項式）の $y^{n-2j}$ の係数は、$U_n = t^n$ の有理式として
\[
c_{2j}(n) = t^j\, \frac{(q/t; t^2)_j}{(t^2; t^2)_j}\cdot\frac{(1 - U_n)(1 + U_n t^{-2j})\prod_{i=1}^{j-1}(1 - U_n^2 t^{-2i})}{\prod_{i=1}^{j}(1 - qU_n^2 t^{-2i-1})}
\]
と書ける。
% [S4]
これは予想で、有理点 4 つで $j \le 5$、$n \le 30$ について漸化式と一致することを確かめた。
virtual 正規化 $\tilde Z = y^N z^N Z$ は $y, z$ の多項式で、$y^a z^b$ の係数は $N$ によらない有限和になる。したがって $U$ を形式変数のまま扱える（$U = t^N$、$N = 3, \dots, 8$ の直接計算と $a + b \le 8$ で全係数一致）。

\subsection{unknot}
% [H1]
\begin{proposition}
unknot では $f_{10} = (1 - U)(U + t)/(U(t-1))$、$f_{a0} = 0$（$a \ge 2$）である。$\nu$ で書くと
\[
f_{10} = -1 - \frac{\nu - \nu^{-1}}{t^{1/2} - t^{-1/2}}
\]
で、Mariño 型の unknot $(\nu - \nu^{-1})/z$ とクロスキャップの $-1$（Mariño の $g_{[1]} = 1$）の和である。積表示は $\prod_j (y_j^{\pm}; t)_\infty / (U y_j^{\pm}; t)_\infty$ で、$\tfrac12$ も $\sum_{d\ \mathrm{odd}}$ も要らない。
\end{proposition}
\begin{proof}
双対 Cauchy 核の期待値は $\prod_j\prod_{k=1}^{N} t^{-(N-k)}(1 - t^{N-k}y_j^{\pm})$ で、$\log$ をとればよい。
\end{proof}

\subsection{Hopf 絡み目}
二つの成分に probe を一つずつ付け、$f_{ab} = \sum_{d | (a,b)} (\mu(d)/d)\, L_{a/d, b/d}(q^d, t^d, U^d)$ と置く。
% [H2][H3]
\begin{compresult}\label{cr:hopf}
$a, b \le 4$ で次が成り立つ。
\begin{enumerate}
\item $f_{ab} \in \dfrac{1 - q}{1 - t}\,\mathbb{Z}[q^\pm, t^\pm, U^\pm]$。$(a, b)$ がともに偶数のときは分母が $1 - t^2$ になる。
\item 成分 1 の probe を無限個にした Schur 基底の $f_{R,b}$（$|R| \le 4$、$b \le 2$）でも同じ整数性が成り立つ。素朴な Macdonald 基底では $(1 - qt)$ などの余分な分母が出る。
\item $f_{11} = -f_{10}\, z_q(\tnu - \tnu^{-1})$、$f_{1b} = f_{11}\,\chi_{b-1}(U/t)$（$\chi_m(x) = x^m + x^{m-2} + \dots + x^{-m}$）。$t = q$ では Mariño の $\hat g_{[1],[1]} = z(\nu - \nu^{-1})$ に、符号と unknot の因子を除いて一致する。
\end{enumerate}
\end{compresult}
refined では、Kähler 変数が $\nu \to \tnu = \nu (q/t)^{1/2}$ とずれる。これは AM の Kähler 変数の量子シフト（§\ref{sec:index}）の開弦版と見なせる。

\subsection{クロスキャップ部分}
被覆には、U$(N)$ の refined CS の composite 表現での Hopf 絡み目を使う。$\tfrac12 f^{\cov}_{10} = -(\nu - \nu^{-1})/z_t$、$\tfrac12 f^{\cov}_{11} = (\nu - \nu^{-1})(\tnu - \tnu^{-1})z_q/z_t$ である。
% [H4][H5]
\begin{compresult}\label{cr:hopfg}
\begin{enumerate}
\item $g_{1b} = z_q(\tnu - \tnu^{-1})\,\chi_{b-1}(\tnu/q^{1/2})$（$b \le 4$）。$t$ は $\tnu$ を通してしか入らない。
\item $\hat g_{R,[1]}$（$|R| \le 4$）は整数 Laurent である。hook $R = (n-d, 1^d)$ では $(-1)^d z_q(\tnu - \tnu^{-1})\tnu^{2d-n+1}$、hook でない $R$ では 0 になる。
$t = q$ では Mariño の $\hat g_{[2],[1]} = z(\nu^2 - 1)$, $\hat g_{[1,1],[1]} = z(\nu^{-2} - 1)$ に、$R \leftrightarrow R^T$ と全体の符号 $-1$ で一致する。
\end{enumerate}
\end{compresult}

\subsection{$E$ の決定と Adams}
probe を $x$ 側と $1/x$ 側に分けたセクター別の被覆から、固定セクター（$\tilde I$ で自分に移るセクター）の状態の和 $P$ を作る。
$E = P - 2Q$（$Q$ は $\tilde I$ の固有値 $-1$ の状態の和、各単項式で $0 \le Q_m/P_m \le 1$）と書き、次の 2 条件を整数計画で課す。
(i) $t = -1$ で $E = 2(1+t)g$。
(ii) $t = q$ で $C$ が Mariño の値。
% [H6][H7]
\begin{compresult}\label{cr:E22}
\begin{align*}
E_{22} &= -z_q (z_q^+)^2(\tnu - \tnu^{-1})(t q^{-1/2}\tnu + q^{1/2}\tnu^{-1})\quad(\text{一意}),\\
E_{24} &= E_{22}\,(\tnu^2/q + q/\tnu^2)\quad(\text{一意}).
\end{align*}
$(4,4)$ は解が存在するが一意でない（固定セクターの単項式 180 個のうち 164 個が自由）。
成分 1 を多 probe にすると、$\tilde I$ は「固定セクターの入れ替え＋Adams」として働き、$E_{R,S} = \sum_{\rho, \sigma\vdash 2} c^R_{2;\rho}c^S_{2;\sigma}\,\varepsilon_{\rho\sigma}$ が成り立つ。確かめた範囲は、$|R| \le 4$, $|S| \le 2$ の 10 セクターと、$|R|, |S| \le 4$ の $t = -1$ での値 25 組である。
\end{compresult}
$P_{22} = z_q^3(\tnu - \tnu^{-1})(tq^{-1/2}\tnu - q^{1/2}\tnu^{-1})$ と比べると、$E_{22}$ は差の因子 $z_q^2$ を和の因子 $-(z_q^+)^2$ に替え、第 2 因子の符号を反転したものである。
これは $\tilde I$ の捻ったトレースの形で、§\ref{sec:positivity} で結び目に当てる。

% [H8]
\begin{compresult}\label{cr:hopfC}
成分 1 を多 probe（$|R| \le 4$）、成分 2 を 2 箱（$|S| \le 2$）にした $\hat C_{R,S} = M^{-1}(q^{-1})\,(g - \tfrac12E/(1+t))\,M^{-1}(q^{-1})^T$ を、$t = q$ の部分（$\tnu$ で表す）と refined だけの部分に分ける。
\begin{enumerate}
\item refined だけの部分は、10 成分すべてで $(\tnu^2 - q)(q - t)$ で割り切れる。太字では $\bt^2(\ba^2 - \bq^2)\cdot\bq^2(\bt^2 - 1)$ である。
\item $t = q$ の部分は 10/10 で、$\tnu$ の冪ごとに符号がそろう。
\item refined だけの部分の商は 8/10 でそろう。転置の対 $\hat C_{(3,1),(1,1)}$, $\hat C_{(2,1,1),(2)}$ では、$\tnu^2$ のブロックに $q^{1/2}(1-t)$ が入ってそろわない。
\end{enumerate}
\end{compresult}
